\subsection{THETA3-C29: actual horizontal boundary special divisors realize the elliptic Poincar\'e kernels}\label{C29}

There is a concrete integral special-cycle realization of the
normalized kernels in C23 in the CM elliptic-power sector.
It uses labelled horizontal boundary special divisors on a
larger unitary Shimura variety. The resulting identity is an
identity of line bundles and actual rational sections over
the integral base. It is not an identification with a
full-rank vertical local intersection.

\paragraph{The precise boundary datum.}
Let $k$ be an imaginary quadratic field of odd discriminant,
and let $\mathcal E$ be a CM elliptic curve with its
$\mathcal O_k$ action and principal polarization over a
base $S\to\mathcal M_{(1,0)}$. One may take
$k=\Q(\sqrt{-7})$. Fix a positive definite self-dual
Hermitian $\mathcal O_k$-lattice $L_0$ of rank $r\geq2$,
and put
\[
                  \mathcal B=\mathcal E\otimes_{\mathcal O_k}L_0.
\]
For a nonzero vector $x\in L_0$, write
$j_x:\mathcal B\to\mathcal E$ for the homomorphism
given by its Hermitian linear functional, with the convention
of the chosen polarization. The additivity relation is
$j_{x+y}=j_x+j_y$. Put $Z(x)=j_x^{-1}(0)$.

These are actual unitary boundary data. Choose the unitary
space with lattice $L_0\perp\mathbb H$, where the
hyperbolic plane has Gram matrix
$\left(\begin{smallmatrix}0&1\\1&0\end{smallmatrix}\right)$.
It has signature $(r+1,1)$ and a rational cusp. In the
notation of Bruinier--Howard--Kudla--Rapoport--Yang,
\emph{Modularity of generating series of divisors on unitary
Shimura varieties}, corrected 2017 author version,
Proposition~3.4.3, pp.~46--48, its boundary abelian scheme
is $B_\Phi=\mathcal E\otimes L_0$ after the stated
CM-base identification. Proposition~3.6.1, pp.~52--53,
identifies each labelled $Z_\Phi(x)$ with $Z(x)$ and
proves that it is an effective Cartier divisor flat over
the CM base. Theorem~3.7.1(4),(8), pp.~53--55,
identifies their sum at a fixed norm with the formal
boundary description of the closure of the corresponding
interior Kudla--Rapoport divisor.

We work with these labelled boundary components. A global
unlabelled fixed-norm divisor sums several such components;
it is not silently replaced by a chosen summand. The formal
boundary chart, and the finite stabilizer quotient when
present, remain part of this comparison. All statements
below may be pulled to a neat level and to the chosen
base $S$.

\paragraph{The normalized line bundle and its section.}
Let $\omega=\operatorname{Lie}(\mathcal E)^{-1}$ and let
$\pi:\mathcal E\to S$. With the principal self-duality,
let $\mathcal P$ be the normalized Poincar\'e bundle on
$\mathcal E\times_S\mathcal E$, and put
\[
       \mathcal Q=\mathcal O_{\mathcal E}(0)\otimes\pi^*\omega.
\]
This is rigidified along zero. The theorem of the square,
in its rigidified form, gives the canonical identity
\[
 \mathcal P\simeq m^*\mathcal Q\otimes
          \operatorname{pr}_1^*\mathcal Q^{-1}\otimes
          \operatorname{pr}_2^*\mathcal Q^{-1}.
 \tag{C29a}
\]
This is the same biextension normalization as C23;
SGA~7~I, VII~2.9.5, p.~175, and VIII~3.1--3.2,
pp.~254--255, identify the rigidified correspondence.

There is a specified integral section $s$ of
$\mathcal Q^{24}$ whose divisor is exactly $24(0)$.
The precise input is the above author version,
Proposition~5.1.4, pp.~69--70: its section is
$\Theta^{24}$ of $J_{0,12}=\mathcal Q^{24}$,
where $\Theta=i\vartheta_1/\eta$ in the stated analytic
convention. That proposition also provides the nowhere
vanishing section of $\omega^{12}$ fixing the associated
Hodge-line normalization. In particular there is no
unrecorded vertical divisor in $s$.

Let $x,y,x+y$ all be nonzero, and set
\[
        \mathcal P_{x,y}=(j_x,j_y)^*\mathcal P.
\]
Then the following rational section is defined on the
dense complement of the indicated divisors:
\[
  s_{x,y}=\frac{j_{x+y}^*s}{j_x^*s\,j_y^*s}
                 \quad\hbox{of }\mathcal P_{x,y}^{24}.
 \tag{C29b}
\]
The exact integral divisor identity is
\[
 \operatorname{div}(s_{x,y})=
       24\bigl(Z(x+y)-Z(x)-Z(y)\bigr),
 \qquad
 \mathcal P_{x,y}^{24}\simeq
       \mathcal O\bigl(24Z(x+y)-24Z(x)-24Z(y)\bigr).
 \tag{C29c}
\]
Consequently
\[
       c_1(\mathcal P_{x,y})
                  =[Z(x+y)]-[Z(x)]-[Z(y)]
                         \quad\hbox{in }\CH^1(\mathcal B)_\Q.
 \tag{C29d}
\]

\textbf{Proof.}
Pull (C29a) back by $(j_x,j_y)$ and take its twenty-fourth
power. It identifies the line of (C29b) with
$\mathcal P_{x,y}^{24}$, with no unspecified scalar.
Each $j_v^*s$ is nonzero as a rational section because
$j_v$ is flat and surjective for $v\ne0$, as used in
the cited proof of Proposition~3.6.1. Its divisor is
$24Z(v)$. Divisors add under tensor products and subtract
under division of nonzero rational sections. This proves
both equalities in (C29c), over the integral base itself.
Taking first Chern classes and dividing by $24$ proves
(C29d). No proper intersection or reducedness of the
individual kernels was required.\hfill$\square$

The power $24$ is deliberately retained: it removes the
theta multiplier and the torsion Hodge-line ambiguity
using the specified integral Jacobi section. At residue
characteristics two or three, dividing a cycle class by
$24$ has a valuation cost. The integral line-bundle and
section statement (C29c) itself has not inverted that
integer. If $x+y=0$, one may replace $y$ by $2y$:
then $x,y, x+2y$ are nonzero, and biextension bilinearity
gives $\mathcal P_{x,2y}=\mathcal P_{x,y}^2$.
Thus the same construction realizes $48c_1(\mathcal P_{x,y})$
as an actual integral combination of horizontal divisors.
If either $x$ or $y$ is zero, the normalized kernel is
trivial and no divisor formula is needed.

\paragraph{An actual C23 kernel, with every label visible.}
Take $L_0=\mathcal O_k^2$ with its standard positive
Hermitian form, so $\mathcal B=\mathcal E^2$. For
$a\in\mathcal O_k\setminus\{0\}$ choose $x,y$ so that
\[
       j_x(P,Q)=[a]P,\qquad j_y(P,Q)=Q.
\]
If the chosen Hermitian convention is linear in the first
entry, $x$ uses the coordinate $\bar a$; this is just the
identity $j_x(P,Q)=\langle(P,Q),x\rangle$.
Their norms are $N_{k/\Q}(a)$, $1$, and
$N_{k/\Q}(a)+1$ for $x,y,x+y$, respectively.
The line $\mathcal P_{x,y}$ is exactly
\[
                     ([a]\operatorname{pr}_1,
                           \operatorname{pr}_2)^*\mathcal P.
 \tag{C29e}
\]
Under the principal polarization this is the normalized
kernel of $h_a=\lambda[a]$ in C23. Its divisor expression
uses precisely the three labelled horizontal boundary
special divisors of these norms. The ambient unitary
variety has signature $(3,1)$; its boundary is a surface.
This is a specified realization, not just an equality of
names for two kinds of divisor.

For powers of $\mathcal E$, an
$\mathcal O_k$-linear homomorphism is a matrix of
endomorphisms $[a]$. Biextension additivity expresses
its normalized Poincar\'e kernel as the tensor product
of the coordinate kernels (C29e), with inverses for
negative coefficients. Apply (C29c) to each nonzero
entry. The selected embedding projectors in C23 are
finite algebraic linear combinations of these classes,
so the same formulas realize those pure kernels with
their exact coefficients. If the auxiliary CM factors
are merely isogenous to such powers, one must specify
the isogenies and retain their pullback and denominator
factors; no integral inverse isogeny is presumed.

\paragraph{Substitution into the actual integral A6 cycle.}
Suppose the auxiliary kernels in C23 belong to this
specified elliptic-power sector, with the required
isogeny factors already fixed. For each term write the
identity (C29c), or its $48$-multiple in the exceptional
zero-sum case, on its appropriate boundary abelian
scheme. Pull the resulting line-bundle identities back
by the actual maps $g_i$ of C23. Substitute their first
Chern identities into its displayed formula for
$\mathcal C$ and expand the finite product.
This expresses the actual cycle $\mathcal C$ as a
finite algebraic-coefficient combination of proper
pushforwards of intersections of these pulled-back
horizontal boundary special-divisor classes.

This assertion is an exact identity in the same Chow
group as C23, because first Chern operators are additive
under tensor product and commute with the indicated
pullbacks and proper pushforward. It is not a new
nonvanishing argument: the nonzero realization is the
one already proved in A6/C23. It now has a specified
special-divisor expression. The factors $24$ and $48$,
the original coefficient field, and every isogeny
denominator remain in the expression.

There is a corresponding statement at the level of
perfect complexes for each effective term. An effective
Cartier divisor $Z(v)$ has the two-term resolution
\[
          [\mathcal O(-Z(v))\xrightarrow{s_v}\mathcal O].
\]
Under an arbitrary comparison map, its derived pullback
is the pulled-back two-term complex. Tensoring several
such complexes retains all their higher Tor groups.
If the pulled-back equations form a regular sequence,
this complex reduces to the usual intersection quotient.
If not, replacing it by the ordinary scheme-theoretic
intersection would lose information. The equality of
first Chern operations above does not assert an equality
between one effective divisor and a signed divisor, or
between their individual structure sheaves.

\paragraph{Why a signed expression is necessary for a pure kernel.}
Let $B_1,B_2$ be positive-dimensional complex abelian
varieties, and let a nonzero class $\gamma$ lie in the
mixed summand $H^1(B_1)\otimes H^1(B_2)$ of $H^2$.
Choose ample classes $\theta_i$ on the factors and put
$g_i=\dim B_i$, $\theta=\operatorname{pr}_1^*\theta_1+
\operatorname{pr}_2^*\theta_2$. Then
\[
       \int_{B_1\times B_2}\gamma\,\theta^{g_1+g_2-1}=0.
 \tag{C29f}
\]
Every summand has odd total cohomological degree on
each factor, whereas a nonzero integral over a factor
requires degree $2g_i$. Thus each term vanishes.
A nonzero effective divisor has positive degree for
an ample polarization: a sufficiently high very ample
power computes that degree by a positive projective
intersection number. Consequently a nonzero pure
mixed kernel class cannot be represented by one
effective divisor. The difference in (C29d) is
therefore essential; no single effective KR divisor
can replace a nonzero pure kernel on this product.

\paragraph{The exact advance and the remaining interior comparison.}
This proves an actual horizontal integral special-cycle
realization for the elliptic-power kernels, on the
specified boundary of a larger unitary integral model.
The fixed-norm sums are the boundary limits of
interior special divisors by the cited toroidal theorem.
The original compact one-active-place ball has no such
rational boundary. Nor does a general CM factor of
dimension greater than one become an elliptic power
merely by calling its polarization a Poincar\'e bundle.
Thus the remaining general comparison is a
correspondence from the chosen compact 032 model to
this formal boundary or to a suitable higher-CM
boundary model, or a direct interior PEL identification.
The construction here supplies the actual horizontal
kernels for its stated sector and preserves the
separate vertical correction in the arithmetic theory.

\framework{The line identities are rigidified divisorial
correspondences in SGA~7~I, VII~2.9.5, p.~175, and
VIII~3.1--3.2, pp.~254--255. Their substitution into
an image cycle is the correspondence construction of
cote~44, scans~7--8. The proper pushforward and
specialization operations use C17's exact SGA/EGA
locators. The signed formula and the derived
pullback distinguish algebraic cycle classes from
local arithmetic lengths.}
